% ============================================================
\chapter{Flot de Ricci~: Introduction}
\label{ch:ricci_flow}
% ============================================================

\section{Motivation et d\'efinition}

Le flot de Ricci, introduit par R.~Hamilton en 1982, est un outil fondamental
en g\'eom\'etrie diff\'erentielle. L'id\'ee est de d\'eformer une m\'etrique riemannienne
dans la direction de sa courbure de Ricci, lissant ainsi la g\'eom\'etrie au cours
du temps.

\begin{intuition}
Le flot de Ricci est l'analogue g\'eom\'etrique de l'\'equation de la chaleur~: tout
comme la chaleur diffuse pour \'egaliser la temp\'erature, le flot de Ricci
\og{}diffuse\fg la courbure pour la rendre plus uniforme. Les r\'egions de forte
courbure positive se contractent, tandis que les r\'egions de courbure n\'egative
s'\'etendent.
\end{intuition}

\begin{definition}[Flot de Ricci]
\index{flot de Ricci}
Soit $(M, g_0)$ une vari\'et\'e riemannienne compacte. Le \emph{flot de Ricci} est
l'\'equation d'\'evolution~:
\[
    \frac{\partial g}{\partial t} = -2\operatorname{Ric}(g), \qquad g(0) = g_0.
\]
\end{definition}

\begin{remarque}
Le facteur $-2$ est une convention. Le signe n\'egatif assure que les r\'egions de
courbure de Ricci positive se contractent (ce qui est physiquement attendu).
L'\'equation est un syst\`eme quasi-lin\'eaire parabolique pour la m\'etrique $g_{ij}$.
\end{remarque}

\begin{definition}[Flot de Ricci normalis\'e]
\index{flot de Ricci normalis\'e}
Le \emph{flot de Ricci normalis\'e} pr\'eserve le volume~:
\[
    \frac{\partial g}{\partial t} = -2\operatorname{Ric}(g) + \frac{2}{n}\bar{R}\, g,
\]
o\`u $\bar{R} = \frac{\int_M R\, dV_g}{\int_M dV_g}$ est la courbure scalaire moyenne
et $n = \dim M$.
\end{definition}

\section{Existence et unicit\'e \`a court terme}

\begin{theoreme}[Hamilton, 1982]
\label{thm:hamilton_existence}
\index{th\'eor\`eme d'existence de Hamilton}
Soit $(M, g_0)$ une vari\'et\'e riemannienne compacte. Il existe $T > 0$ et une
unique famille lisse de m\'etriques $(g(t))_{t \in [0, T)}$ satisfaisant le flot de Ricci
avec condition initiale $g(0) = g_0$.
\end{theoreme}

\begin{remarque}
La preuve originale de Hamilton utilise le th\'eor\`eme de Nash-Moser. DeTurck (1983)
a ensuite donn\'e une preuve simplifi\'ee en introduisant une modification du flot
qui le rend strictement parabolique.
\end{remarque}

\begin{definition}[Astuce de DeTurck]
\index{astuce de DeTurck}
Soit $\tilde{g}$ une m\'etrique de r\'ef\'erence fixe sur $M$. Le \emph{flot de
Ricci-DeTurck} est~:
\[
    \frac{\partial g}{\partial t} = -2\operatorname{Ric}(g)
    + \mathcal{L}_{W(g, \tilde{g})} g,
\]
o\`u $W$ est un champ de vecteurs d\'ependant de $g$ et $\tilde{g}$ (d\'efini via
la diff\'erence des symboles de Christoffel). Ce flot est strictement parabolique.
\end{definition}

\begin{proposition}
Le flot de Ricci-DeTurck est \'equivalent au flot de Ricci modulo diff\'eomorphismes~:
si $g(t)$ est une solution du flot de Ricci-DeTurck, il existe une famille de
diff\'eomorphismes $\varphi_t : M \to M$ telle que $\varphi_t^* g(t)$ soit une solution
du flot de Ricci.
\end{proposition}

\section{\'Evolution de la courbure}

\begin{theoreme}[\'Evolution de la courbure scalaire]
\label{thm:evolution_R}
Sous le flot de Ricci, la courbure scalaire \'evolue selon~:
\[
    \frac{\partial R}{\partial t} = \Delta R + 2\abs{\operatorname{Ric}}^2.
\]
\end{theoreme}

\begin{proof}
En coordonn\'ees locales, $R = g^{ij} R_{ij}$. En d\'erivant~:
\[
    \frac{\partial R}{\partial t} = g^{ij} \frac{\partial R_{ij}}{\partial t}
    + \frac{\partial g^{ij}}{\partial t} R_{ij}.
\]
Comme $\frac{\partial g^{ij}}{\partial t} = 2R^{ij}$ et en utilisant la formule
d'\'evolution du tenseur de Ricci (obtenue par un calcul tensoriel direct)~:
\[
    \frac{\partial R_{ij}}{\partial t} = \Delta R_{ij} + 2R_{ikjl}R^{kl}
    - 2R_{ik}R_j^{\;k},
\]
on obtient apr\`es contraction le r\'esultat annonc\'e.
\end{proof}

\begin{corollaire}[Principe du maximum pour $R$]
Si la courbure scalaire initiale v\'erifie $R(g_0) \geq R_{\min}$, alors
pour tout $t \in [0, T)$~:
\[
    R(g(t)) \geq \frac{R_{\min}}{1 - \frac{2R_{\min}}{n} t}.
\]
En particulier, si $R_{\min} > 0$, le flot d\'eveloppe une singularit\'e en temps fini
$T \leq \frac{n}{2R_{\min}}$.
\end{corollaire}

\begin{theoreme}[\'Evolution du tenseur de Ricci]
Sous le flot de Ricci~:
\[
    \frac{\partial R_{ij}}{\partial t} = \Delta R_{ij}
    + 2R_{ikjl}R^{kl} - 2R_{ik}R_j^{\;k}.
\]
\end{theoreme}

\begin{theoreme}[\'Evolution du tenseur de Riemann]
Sous le flot de Ricci~:
\[
    \frac{\partial R_{ijkl}}{\partial t} = \Delta R_{ijkl}
    + 2(B_{ijkl} - B_{ijlk} + B_{ikjl} - B_{iljk}),
\]
o\`u $B_{ijkl} = R_{ipjq}R_{kplq}\, g^{pp'} g^{qq'}$ (on somme sur les indices
r\'ep\'et\'es avec la convention d'Einstein).
\end{theoreme}

\section{Principe du maximum de Hamilton}

\begin{theoreme}[Principe du maximum scalaire]
\label{thm:maximum_scalaire}
\index{principe du maximum}
Soit $u : M \times [0, T) \to \R$ une solution de~:
\[
    \frac{\partial u}{\partial t} \geq \Delta_{g(t)} u + F(u),
\]
o\`u $g(t)$ est une solution du flot de Ricci et $F$ est localement lipschitzienne.
Si $u(x, 0) \geq c$ pour tout $x \in M$, alors $u(x, t) \geq \phi(t)$ o\`u $\phi$
est la solution de l'EDO $\phi' = F(\phi)$, $\phi(0) = c$.
\end{theoreme}

\begin{proof}[Esquisse de preuve]
On utilise le principe de comparaison parabolique. Soit $x_0(t)$ un point o\`u
$u(\cdot, t)$ atteint son minimum spatial. En ce point, $\Delta_{g(t)} u \geq 0$
(car c'est un minimum), donc~:
\[
    \frac{\partial u}{\partial t}\bigg|_{x_0(t)} \geq F(u(x_0(t), t)).
\]
On en d\'eduit (par un argument de r\'egularisation, car $x_0(t)$ peut sauter)
que $u_{\min}(t) = \min_M u(\cdot, t)$ satisfait $\frac{d}{dt} u_{\min} \geq F(u_{\min})$
au sens des barri\`eres. Par comparaison avec l'EDO $\phi' = F(\phi)$, $\phi(0) = c$,
on obtient $u_{\min}(t) \geq \phi(t)$, d'o\`u le r\'esultat.
L'application typique est $u = R$ (courbure scalaire), $F(u) = \frac{2}{n}u^2$,
donnant la borne $R \geq \frac{R_{\min}}{1 - \frac{2R_{\min}}{n}t}$.
\end{proof}

\begin{attention}[title=\textbf{La m\'etrique d\'epend du temps}]
Le laplacien $\Delta_{g(t)}$ d\'epend de la m\'etrique qui \'evolue. Le principe du
maximum n\'ecessite donc un traitement soigneux, car les boules m\'etriques et les
normes changent au cours du temps.
\end{attention}

\begin{theoreme}[Principe du maximum tensoriel de Hamilton]
\label{thm:maximum_tensoriel}
Soit $T_{ij}$ un tenseur sym\'etrique \'evoluant par~:
\[
    \frac{\partial T_{ij}}{\partial t} = \Delta T_{ij} + N_{ij}(T, g),
\]
o\`u $N_{ij}$ satisfait~: si $T_{ij} \geq 0$ et $T_{ij}v^j = 0$ en un point,
alors $N_{ij} v^i v^j \geq 0$. Si $T_{ij}(0) \geq 0$, alors $T_{ij}(t) \geq 0$
pour tout $t \in [0, T)$.
\end{theoreme}

\begin{proof}[Esquisse de preuve]
On consid\`ere la plus petite valeur propre $\lambda_{\min}(t)$ de $T_{ij}$ sur $M$.
Si $\lambda_{\min}(t_0) = 0$ en un point $x_0$ avec vecteur propre $v$, alors
$T_{ij}v^j = 0$ en ce point, donc $N_{ij}v^i v^j \geq 0$ par hypoth\`ese. Un
argument de principe du maximum (appliqu\'e \`a la quantit\'e $T_{ij}v^i v^j$ le long
de courbes de vecteurs parall\`eles) montre que $\lambda_{\min}$ ne peut franchir
z\'ero par valeurs d\'ecroissantes. La difficult\'e technique r\'eside dans le fait
que le vecteur propre d\'epend du point~; on la contourne en travaillant avec le
fibr\'e des rep\`eres orthonorm\'es.
Voir~\textsc{Hamilton}, \emph{Three-manifolds with positive Ricci curvature},
J.~Differential Geom.~17 (1982), Theorem~9.1.
\end{proof}

\begin{corollaire}[Pr\'eservation de la positivit\'e de Ricci]
Si $\operatorname{Ric}(g_0) \geq 0$, alors $\operatorname{Ric}(g(t)) \geq 0$
pour tout $t$ dans l'intervalle d'existence du flot.
\end{corollaire}

\section{Exemples fondamentaux}

\begin{exemple}[Sph\`ere $S^n$]
\index{flot de Ricci!sph\`ere}
Sur la sph\`ere $(S^n, g_0)$ avec $g_0$ la m\'etrique ronde de courbure sectionnelle $1$,
on a $\operatorname{Ric}(g_0) = (n-1)g_0$. Le flot de Ricci a pour solution~:
\[
    g(t) = (1 - 2(n-1)t)\, g_0.
\]
La m\'etrique se contracte uniform\'ement et s'annule en $T = \frac{1}{2(n-1)}$.
La sph\`ere se contracte en un point en temps fini.
\end{exemple}

\begin{exemple}[Espace hyperbolique]
\index{flot de Ricci!espace hyperbolique}
Sur $(\mathbb{H}^n, g_0)$ avec $\operatorname{Ric}(g_0) = -(n-1)g_0$~:
\[
    g(t) = (1 + 2(n-1)t)\, g_0.
\]
La m\'etrique s'expand uniform\'ement. Le flot existe pour tout $t > 0$.
\end{exemple}

\begin{exemple}[Solitons de Ricci]
\index{soliton de Ricci}
Un \emph{soliton de Ricci} est une solution du flot qui n'\'evolue que par
diff\'eomorphismes et changement d'\'echelle~:
\[
    \operatorname{Ric}(g) + \mathcal{L}_X g + \lambda g = 0,
\]
o\`u $X$ est un champ de vecteurs et $\lambda \in \R$. Le soliton est dit
r\'etr\'ecissant ($\lambda > 0$), stationnaire ($\lambda = 0$), ou expansif ($\lambda < 0$).
\end{exemple}

\begin{exemple}[Flot sur les surfaces]
\index{flot de Ricci!surfaces}
En dimension $2$, $\operatorname{Ric} = \frac{R}{2}g$ et le flot de Ricci s'\'ecrit~:
\[
    \frac{\partial g}{\partial t} = -Rg.
\]
Le flot de Ricci normalis\'e est $\frac{\partial g}{\partial t} = (\bar{R} - R)g$.
Hamilton (1988) et Chow (1991) ont montr\'e que sur toute surface compacte, le flot
normalis\'e converge vers une m\'etrique \`a courbure constante.
\end{exemple}

\section{Convergence en dimension 3}

\begin{theoreme}[Hamilton, 1982]
\label{thm:hamilton_3d}
Soit $(M^3, g_0)$ une vari\'et\'e riemannienne compacte de dimension $3$ \`a courbure
de Ricci strictement positive. Le flot de Ricci normalis\'e converge vers une m\'etrique
\`a courbure sectionnelle constante positive. En particulier, $M$ est diff\'eomorphe
\`a un quotient $S^3/\Gamma$.
\end{theoreme}

\begin{remarque}
Ce r\'esultat fondateur a motiv\'e le programme de Hamilton-Perelman, aboutissant
\`a la preuve de la conjecture de Poincar\'e par Perelman (2002--2003) en dimension $3$.
\end{remarque}

\begin{formules}[title=\textbf{\'Equations d'\'evolution sous le flot de Ricci}]
\begin{align*}
    \frac{\partial g_{ij}}{\partial t} &= -2R_{ij} \\[6pt]
    \frac{\partial}{\partial t}\Gamma_{ij}^k &= -g^{kl}(\nabla_i R_{jl}
        + \nabla_j R_{il} - \nabla_l R_{ij}) \\[6pt]
    \frac{\partial R}{\partial t} &= \Delta R + 2\abs{\operatorname{Ric}}^2 \\[6pt]
    \frac{\partial R_{ij}}{\partial t} &= \Delta R_{ij} + 2R_{ikjl}R^{kl}
        - 2R_{ik}R_j^{\;k} \\[6pt]
    \frac{\partial}{\partial t}dV_g &= -R\, dV_g
\end{align*}
\end{formules}

\section{Exercices}

\begin{exercice}
V\'erifier que $g(t) = (1 - 2(n-1)t) g_0$ est bien une solution du flot de Ricci
sur $(S^n, g_0)$. Calculer le temps d'extinction et le diam\`etre $\operatorname{diam}(S^n, g(t))$
en fonction de $t$.
\end{exercice}

\begin{exercice}
Montrer que sous le flot de Ricci, le volume \'evolue selon~:
$\frac{d}{dt}\operatorname{Vol}(M, g(t)) = -\int_M R\, dV_g$.
En d\'eduire que le volume d\'ecro\^it si $R > 0$.
\end{exercice}

\begin{exercice}
Soit $(M^2, g_0)$ une surface compacte. Montrer que le flot de Ricci normalis\'e
pr\'eserve l'aire totale et que la caract\'eristique d'Euler-Poincar\'e $\chi(M)$
d\'etermine le signe de $\bar{R}$ via Gauss-Bonnet.
\end{exercice}

\begin{exercice}[Soliton cigare de Hamilton]
Montrer que la m\'etrique sur $\R^2$~:
\[
    g = \frac{dx^2 + dy^2}{1 + x^2 + y^2}
\]
est un soliton de Ricci stationnaire (dit \og{}soliton cigare\fg).
Calculer sa courbure de Gauss.
\end{exercice}

\begin{exercice}
Soit $g(t) = \varphi(t) g_0$ une solution du flot de Ricci avec $\varphi(0) = 1$.
Montrer que cela n'est possible que si $(M, g_0)$ est un soliton d'Einstein~:
$\operatorname{Ric}(g_0) = \lambda g_0$. D\'eterminer $\varphi(t)$ en fonction de
$\lambda$.
\end{exercice}

\begin{exercice}
En utilisant le principe du maximum, montrer que si $R(g_0) \geq R_{\min} > 0$
sur une vari\'et\'e compacte, alors $R(g(t)) \to +\infty$ lorsque $t$ tend vers
le temps maximal d'existence.
\end{exercice}
